\section{Elementary certificates for the numerical inequalities}
\label{cert:arithmetic}
This appendix gives rational enclosures for the numerical comparisons
used in the proof. All finite decimals denote exact rationals. The
formulas below and the tables suffice to reproduce the comparisons;
no numerical optimization or numerical quadrature is needed.

\subsection{Enclosure formulas}
We use
\begin{equation}\label{cert:pi}
 3.14159<\pi<3.14160.
\end{equation}
For example, these bounds follow from
\(\pi=16\arctan(1/5)-4\arctan(1/239)\) by bounding each arctangent
between its alternating-series partial sums through indices 7 and 6.
All square-root enclosures may be obtained using rational numbers
\(l,u\ge0\) with \(l^2\le x\le u^2\).
An explicit rule, if desired, is to set \(D=10^{35}\) and use
\[
 \frac{\lfloor\sqrt{\lfloor xD^2\rfloor}\rfloor}{D}
 \le\sqrt{x}\le
 \frac{\lfloor\sqrt{\lfloor xD^2\rfloor}\rfloor+1}{D}.
\]
For an interval of possible \(x\), apply the lower rule to its lower
endpoint and the upper rule to its upper endpoint. Thus only rational
operations and integer square roots are required.

For \(y\ge0\), let
\[
 E_N(y)=\sum_{j=0}^N\frac{(-y)^j}{j!}.
\]
Taylor's integral remainder gives
\begin{equation}\label{cert:exponential}
 E_{15}(y)\le e^{-y}\le E_{14}(y).
\end{equation}
This conclusion does not require the initial terms of the series to
be decreasing. Integrating the same bounds, set
\[
 J_N(x)=\sum_{j=0}^N
 \frac{(-1)^j x^{2j+1}}{2^j j!(2j+1)},\qquad x\ge0.
\]
Then
\begin{equation}\label{cert:normal}
 \frac12+\frac{J_{15}(x)}{\sqrt{2\pi}}
 \le\Phi(x)\le
 \frac12+\frac{J_{14}(x)}{\sqrt{2\pi}},\qquad
 \Phi(-x)=1-\Phi(x).
\end{equation}
Together with \eqref{cert:pi} and the square-root rule, these are
rational enclosures for both \(\phi\) and \(\Phi\) at every argument
used below. Interval addition, subtraction, multiplication, and
division by intervals avoiding zero preserve the enclosures.

For the cap calculation, put
\[
 Q(x)=\sum_{j=0}^4\frac{\binom{2j}{j}}{4^j(2j+1)}x^{2j+1}.
\]
The binomial series has positive coefficients at most one. Its tail
starting in degree 10 is bounded by \(x^{10}/(1-x^2)\); after
integration this gives, for \(0\le x<1\),
\begin{equation}\label{cert:arcsine}
 Q(x)\le\arcsin x\le Q(x)+\frac{x^{11}}{11(1-x^2)}.
\end{equation}

\subsection{The two rearrangement minima}
Write \(f(x)=(9/4)e^{x^2}\Phi(-x)\) and
\(D_f(x)=2x-\phi(x)/\Phi(-x)\), as in
\eqref{sc:halfspace-ratio}. Substitution into
\eqref{cert:exponential}--\eqref{cert:normal} gives the following
strict enclosures:
\[
\begin{array}{rcl}
 -0.000131&<&D_f(0.6119)<-0.000126,\\
  0.000120&<&D_f(0.6121)< 0.000124,\\
 -0.000430&<&D_f(x)<0.000424
               \quad(0.6119\le x\le0.6121),\\
  0.884383832&<&f(0.612).
\end{array}
\]
For the third row, use the whole rational interval
\([0.6119,0.6121]\) in the elementary interval operations;
\(\phi(x)\) decreases and \(\Phi(-x)\) decreases there, so endpoint
enclosures give both factors. Thus \(|D_f|<0.001\) throughout that
interval. The minimizer lies within \(0.0001\) of \(0.612\), and so
\[
 \min f\ge f(0.612)e^{-10^{-7}}
 >0.884383832(1-10^{-7})>0.8843837.
\]

For \(p,F,H,R\) from \eqref{sc:cap-ratio} and its following
calculation, \eqref{cert:arcsine} gives
\[
\begin{array}{rcl}
 -0.0047010&<&F(0.29)<-0.0046981,\\
  0.0007683&<&F(0.294)<0.0007712,\\
  0.9312526&<&R(0.29),\\
  0.22198&<&p(0.294)-0.294/\pi.
\end{array}
\]
For \(0.29\le x\le0.294\), the last row implies
\[
 H(x)=p(x)-\frac{x\sqrt{1-x^2}}\pi
 \ge p(0.294)-\frac{0.294}\pi>0.
\]
Therefore \(F\) is increasing throughout the interval, and the first
two rows imply \(|F(x)|<0.004702\) there. Consequently
\[
 |R'(x)|<\frac{(9/4)0.004702}{(1-0.294^2)^4}
 <0.015189<0.03.
\]
This proves the derivative bound on the entire interval. The minimum
therefore exceeds \(0.9312526-0.03(0.004)=0.9311326\).

\subsection{Quantiles and the interval probability bounds}
For \(P_0(a)=(0.415+0.15a)a\), Table~\ref{cert:quantiles}
gives a rational \(q_a\) satisfying \(\Phi(q_a)<P_0(a)\), and a
strict upper bound on \(G(q_a)\). Each entry follows by substituting
the displayed rational argument into
\eqref{cert:exponential}--\eqref{cert:normal} and the definition of
\(G\). Since \(G\) is nonincreasing,
\[
 G\bigl(\Phi^{-1}(P_0(r))\bigr)\le G(q_a)
 \qquad(r\ge a).
\]
In each row the listed upper bound is strictly below the value of
\(U\) in Table~\ref{sc:interval-table}.

\begin{table}[ht]
\centering
\begin{tabular}{rrrr}
\hline
\(a\)&\(q_a\)&\(10^6(P_0(a)-\Phi(q_a))>\)&\(G(q_a)<\)\\
\hline
0.05&-2.033&98&0.887728\\
0.14&-1.548&228&0.854397\\
0.24&-1.238&381&0.826606\\
0.33&-1.024&366&0.803691\\
0.41&-0.860&470&0.783737\\
0.48&-0.728&452&0.766006\\
0.54&-0.621&539&0.750456\\
0.59&-0.535&739&0.737150\\
0.64&-0.450&684&0.723262\\
\hline
\end{tabular}
\caption{Rational enclosures certifying the quantile and \(G\) bounds.}
\label{cert:quantiles}
\end{table}

Finally, for each row \((a,b,U)\) of
Table~\ref{sc:interval-table}, replace \(\pi\) by its upper bound in
\(B^2=9/(8\pi)\), enclose \(\sqrt{1-a^2}\) from above, and subtract
the exact rational \(U(0.415+0.15b)^2\). This gives the lower bounds
in that table's last column. The same rational and square-root rules
apply to the finite comparisons in the numerical elimination section.
